\chapter{Construire le modèle}

\section{Pourquoi une table plutôt qu'une formule}
Le chapitre 3 a montré une réponse à trois variables, en cloche selon la
clarté, décroissante selon la saturation, avec des particularités par teinte.
On pourrait chercher une formule à quelques constantes. Chaque hypothèse de
forme risque pourtant de reproduire l'erreur des anciens modèles : imposer un
comportement que Lightroom n'a pas.

Une \textbf{table de réponse} évite ce piège : on découpe l'espace des pixels
en cellules et on stocke, pour chacune, la réponse mesurée. C'est la même
démarche que pour les hautes lumières dans \IPR{} : une référence mesurée,
interpolée, plutôt qu'un modèle paramétrique deviné.

\section{Le découpage}
Trois axes, ceux du pixel neutre :
\begin{itemize}
\item saturation ProPhoto $s$ : 12 classes, de bornes
0 ; 0,05 ; 0,1 ; 0,15 ; 0,2 ; 0,3 ; 0,4 ; \ldots ; 0,9 ; 1 (plus fines près du
gris, où la réponse varie vite) ;
\item clarté Oklab $L$ : 9 classes, de bornes 0 ; 0,2 ; 0,3 ; \ldots ; 0,9 ; 1,2 ;
\item teinte Oklab $h$ : 12 secteurs de $30^\circ$.
\end{itemize}
Soit $12\times9\times12=1\,296$ cellules par valeur mesurée et par grandeur
(gain de chroma, gain de clarté, rotation), pour six valeurs du curseur.

\section{Remplir une cellule}
La valeur d'une cellule est la \textbf{médiane}, sur toutes les scènes, des
$\gamma$ (ou $\lambda$, ou $\Delta h$) de ses pixels. La médiane résiste aux
pixels aberrants (bruit, bord d'un objet) mieux que la moyenne.

Une cellule doit contenir au moins 40 pixels pour être jugée fiable. Sinon,
elle reçoit une valeur de repli :
\begin{itemize}
\item pour les gains : la médiane de sa classe (saturation, clarté), toutes
teintes confondues, puis de sa classe de saturation seule ;
\item pour la rotation : la médiane de sa classe (saturation, teinte), puis de
sa teinte seule, puis zéro.
\end{itemize}
La différence est essentielle. Un gain varie peu d'une teinte à l'autre, une
rotation dépend d'abord de la teinte : les bleus tournent dans un sens, les
oranges dans l'autre. Un premier essai qui repliait aussi la rotation « toutes
teintes confondues » donnait à des cellules bleues vides la rotation des
oranges, majoritaires dans le corpus, donc le mauvais sens.

\resultat{Une valeur de repli doit conserver la variable dont la grandeur
dépend le plus.}

\section{Interpoler entre les cellules}
Un pixel tombe rarement au centre d'une cellule. Prendre la valeur de sa
cellule créerait des marches visibles dans un dégradé de ciel. Le modèle
interpole donc entre les centres des cellules voisines, linéairement sur
chaque axe (interpolation \emph{trilinéaire}). Sur un axe, avec deux centres
$x_0<x_1$ et leurs valeurs $v_0,v_1$ :
\begin{equation}
v(x)=(1-t)\,v_0+t\,v_1,\qquad t=\frac{x-x_0}{x_1-x_0}\in[0,1].
\end{equation}
Au-delà du premier ou du dernier centre, la valeur reste constante. La teinte
est un axe circulaire : le secteur centré sur $345^\circ$ a pour voisin celui
centré sur $15^\circ$.

\section{Entre les valeurs mesurées du curseur}
Les tables existent pour $\pm25$, $\pm50$ et $\pm100$. Pour une valeur
intermédiaire $V$, on interpole linéairement entre les deux valeurs mesurées
qui l'encadrent, du même côté de zéro, avec zéro pour origine (réponse nulle).
L'interpolation se fait sur $\gamma$ et $\lambda$, c'est-à-dire en
logarithme : les gains se combinent par multiplication, et cette interpolation
garantit qu'un gain reste positif.

\section{Monotonie : un curseur cohérent}
Une médiane reste une estimation bruitée. Dans quelques cellules peu
peuplées, elle donnait par exemple un gain plus faible à $+50$ qu'à $+25$. Le
photographe verrait alors la couleur \emph{diminuer} en poussant le curseur.

Le modèle impose donc, cellule par cellule, que $\gamma$ croisse avec la
valeur du curseur. Il utilise la \textbf{régression isotone} : la suite
croissante la plus proche, au sens des moindres carrés, des valeurs mesurées.
L'algorithme \emph{Pool Adjacent Violators}~\cite{ayer1955} la calcule en un
passage : chaque fois que deux valeurs voisines violent l'ordre, on les
remplace par leur moyenne, et on recommence si le bloc obtenu viole l'ordre
avec son voisin.

Exemple : les valeurs $(0{,}162\,;\ 0{,}090\,;\ 0{,}924)$ pour $+25$, $+50$,
$+100$ violent l'ordre entre $+25$ et $+50$. On les remplace par leur moyenne :
$(0{,}126\,;\ 0{,}126\,;\ 0{,}924)$.

La contrainte est appliquée séparément de chaque côté de zéro, qui reste fixe.
Elle a modifié 61 cellules sur 1\,296, sans dégrader la validation croisée.

\section{Les quasi-gris}
Les tables sont construites sur des pixels de chroma supérieure à 0,01. Pour
les quasi-gris, on mesure un \textbf{facteur} $f$ qui multiplie le gain prédit
par la table, par classe de chroma et par valeur du curseur. Il est obtenu par
moindres carrés :
\begin{equation}
f=\frac{\sum_i \gamma_i^{\text{LR}}\,\gamma_i^{\text{table}}}{\sum_i \left(\gamma_i^{\text{table}}\right)^2}.
\end{equation}
À $+100$, il vaut 0,74 près du gris et remonte à 1 vers $C=0{,}016$ : le
renforcement est atténué. À $-100$, il vaut 1,38 : la désaturation est
renforcée. À $\pm25$ et $\pm50$, il reste proche de 1.

Un premier essai calculait un seul facteur par sens du curseur. La valeur
$-100$, aux logarithmes très grands, dominait alors l'ajustement et
désaturait trop les quasi-gris à $-25$ et $-50$ : un facteur par valeur a
corrigé ce biais.

Enfin, sous $C=0{,}01$, le gain de clarté et la rotation s'estompent
proportionnellement à $C/0{,}01$. Un gris exact reste exactement gris, et un
presque-gris ne change pas brusquement de teinte, ce qui aurait révélé le
bruit.

\section{La saturation}
La saturation n'utilise pas de table : sa mesure sur 49 scènes suit
Lightroom à moins de 0,05 de variation de chroma. Elle multiplie la chroma
par $1+S/100$ et conserve $L$ et $h$. Combinée à la vibrance, le gain total de
chroma est $(1+S/100)\,e^{\gamma}$.

\section{Exemple déroulé}
Un bleu de ciel, en RGB linéaire sRGB $(0{,}10\,;\ 0{,}20\,;\ 0{,}45)$, à
vibrance $+100$.
\begin{enumerate}
\item Oklab : $L=0{,}583$, $a=-0{,}017$, $b=-0{,}093$, soit $C=0{,}094$ et
$h=259{,}5^\circ$.
\item ProPhoto linéaire : $(0{,}182\,;\ 0{,}197\,;\ 0{,}415)$. La saturation
vaut donc $s=(0{,}415-0{,}182)/0{,}415$, soit $s=0{,}561$.
\item Position dans les tables : $s$ entre les centres 0,55 et 0,65
($t=0{,}11$) ; $L$ entre 0,55 et 0,65 ($t=0{,}33$) ; $h$ entre les secteurs
centrés sur $255^\circ$ et $285^\circ$ ($t=0{,}15$).
\item Lecture interpolée à $+100$ : $\gamma=0{,}541$ (gain $1{,}72$),
$\lambda=-0{,}043$ (clarté $\times0{,}96$), rotation $-4{,}7^\circ$. Le pixel
n'est pas un quasi-gris ($C>0{,}016$) : facteur 1, pas d'estompage.
\item Résultat : $L'=0{,}558$, $C'=0{,}162$ et $h'=254{,}8^\circ$. De retour en RGB
linéaire sRGB, le pixel vaut $(0{,}015\,;\ 0{,}172\,;\ 0{,}628)$.
\end{enumerate}
Le bleu est plus intense, un peu plus sombre, et a tourné vers le cyan.
